\documentclass[11pt,a4paper]{article}
\usepackage[margin=25mm]{geometry}
\usepackage[T1]{fontenc}
\usepackage{lmodern}
\usepackage{amsmath,amssymb,amsthm}
\usepackage{graphicx}
\usepackage{xcolor}
\usepackage{microtype}
\usepackage[hidelinks]{hyperref}
\usepackage[most]{tcolorbox}

\setlength{\parindent}{0pt}
\setlength{\parskip}{5pt}
\emergencystretch=2em

\newtheorem{theorem}{Theorem}
\newtheorem{lemma}{Lemma}
\newcommand{\C}{\mathbb C}
\newcommand{\D}{\mathbb D}
\newcommand{\T}{\mathbb T}
\newcommand{\Ber}{\mathrm{Ber}}

\title{The relative Berkovich boundary inclusion fails\\
\large A disk-algebra counterexample for arXiv:1410.5893}
\author{}
\date{17 August 2026}

\begin{document}
\maketitle
\vspace{-1.4em}

\begin{tcolorbox}[colback=blue!5,colframe=blue!55!black]
\textbf{Outcome.}
The final question of Leung--Ng has a negative answer, already for commutative
unital complex Banach algebras.  Take
\[
 k=\C,\qquad A=C(\T),\qquad a(\zeta)=\zeta,
 \qquad B=\overline{\C[a]}=A(\D)|_{\T}.
\]
Then
\[
 \mu_{\C}(0)\in
 \partial\sigma^{\Ber}_{B,\C}(a)
 \quad\text{but}\quad
 \mu_{\C}(0)\notin\sigma^{\Ber}_{A,\C}(a),
\]
where the boundary is taken, as in the source, in the full affine analytic
line $\mathbb A^1_{\mathbb Z_1}$.
\end{tcolorbox}

\section{The exact question}

For a unital $k$-Banach algebra $A$, an element $a\in A$, and the closed
unital $k$-Banach subalgebra $B$ generated by $a$, the source asks whether
\[
 \partial\sigma^{\Ber}_{B,k}(a)
 \subseteq \sigma^{\Ber}_{A,k}(a). \tag{1}
\]
The convention established immediately before this residual question is that
the boundary is taken in $\mathbb A^1_{\mathbb Z_1}$, not merely in the slice
of points extending the valuation of $k$.

\begin{figure}[h]
\centering
\includegraphics[width=.92\textwidth]{figures/open_question_crop.png}
\caption{The residual question on source PDF page 25.}
\end{figure}

That ambient-space convention is decisive: even for $k=\C$, the relative
spectrum has transverse valuation directions arbitrarily close to each of its
points.

\section{Every complex-relative spectral point is a boundary point}

Write $\C_1=(\C,|\cdot|)$ and, for $0<\nu<1$, let
$\C_\nu=(\C,|\cdot|^\nu)$.  If $z\in\C$, denote by
$\mu_{\C_\nu}(z)\in\mathbb A^1_{\mathbb Z_1}$ the multiplicative seminorm
\[
 p\longmapsto |p(z)|^\nu \qquad(p\in\mathbb Z[t]).
\]

\begin{lemma}\label{lem:boundary}
Let $E$ be any unital complex Banach algebra and $x\in E$.  Then
\[
 \partial\sigma^{\Ber}_{E,\C_1}(x)
 =\sigma^{\Ber}_{E,\C_1}(x)
\]
as subsets of $\mathbb A^1_{\mathbb Z_1}$.
\end{lemma}

\begin{proof}
The relative spectrum is compact, hence closed, so its boundary is contained
in it.  Conversely, take $\xi$ in the relative spectrum.  By Proposition
4.13(b) of the source,
$\xi=\mu_{\C_1}(z)$ for some $z$ in the usual complex spectrum of $x$.

Choose $\nu_j\uparrow1$ with $\nu_j<1$.  Pointwise on $\mathbb Z[t]$,
\[
 \mu_{\C_{\nu_j}}(z)(p)=|p(z)|^{\nu_j}\longrightarrow |p(z)|
 =\mu_{\C_1}(z)(p),
\]
so $\mu_{\C_{\nu_j}}(z)\to\xi$ in $\mathbb A^1_{\mathbb Z_1}$.
None of these approximating points belongs to a spectrum relative to
$\C_1$.  Indeed, every complete valuation-field extension of $\C_1$ is
isometric on the scalar field and hence assigns the constant polynomial $2$
the value $2$, whereas
\[
 \mu_{\C_{\nu_j}}(z)(2)=2^{\nu_j}\ne2.
\]
Thus every neighborhood of $\xi$ meets the complement of the relative
spectrum, proving that $\xi$ is a boundary point.
\end{proof}

\section{The disk-algebra counterexample}

Let $A=C(\T)$ with the supremum norm and pointwise multiplication, and put
$a(\zeta)=\zeta$.  Let $B$ be the closed unital subalgebra generated by $a$.
By the maximum-modulus theorem, restriction to $\T$ identifies $B$
isometrically with the disk algebra $A(\D)$: it is exactly the uniform closure
of analytic polynomials in the coordinate function.  Thus this is precisely
the generated subalgebra required in (1), with no change of norm.

\begin{theorem}
For this inclusion $B\subset A$,
\[
 \mu_{\C_1}(0)\in\partial\sigma^{\Ber}_{B,\C_1}(a)
 \setminus\sigma^{\Ber}_{A,\C_1}(a).
\]
Consequently, the inclusion in \emph{(1)} is false.
\end{theorem}

\begin{proof}
First compute the ordinary spectra.  In $B$, if $|\lambda|>1$, the Neumann
series
\[
 (\lambda-a)^{-1}=\lambda^{-1}\sum_{n\ge0}(a/\lambda)^n
\]
converges in $B$.  If $|\lambda|\le1$, evaluation at $\lambda$ is a continuous
unital character on $A(\D)$ and annihilates $a-\lambda$, so $a-\lambda$ is not
invertible.  Hence
\[
 \sigma_B^{\C}(a)=\overline{\D}.
\]
In $A=C(\T)$, the function $a-\lambda$ is invertible exactly when it has no
zero on $\T$, and therefore
\[
 \sigma_A^{\C}(a)=\T.
\]

Proposition 4.13(b) of the source identifies the relative Berkovich spectra
with the folded ordinary spectra:
\[
 \sigma^{\Ber}_{B,\C_1}(a)=\mu_{\C_1}(\overline\D),
 \qquad
 \sigma^{\Ber}_{A,\C_1}(a)=\mu_{\C_1}(\T).
\]
Thus $\mu_{\C_1}(0)$ belongs to the first set and not the second.  Lemma
\ref{lem:boundary} says that it is a boundary point of the first set in the
source's ambient topology, which proves the claim.
\end{proof}

\section{Why the counterexample is robust}

The failure is not caused by noncommutativity, an exotic field, a nonisometric
embedding, or a potentially degenerate completed tensor product.  Both
algebras are familiar commutative uniform complex Banach algebras, and $B$ is
literally the inherited-norm closed subalgebra generated by one element.

The mechanism is structural.  Lemma \ref{lem:boundary} shows that, with
boundary taken in the full $\mathbb A^1_{\mathbb Z_1}$, (1) would force
\emph{full} spectral permanence for every complex generated subalgebra, not
only the classical boundary inclusion.  The disk algebra has the standard
spectral hole $\D$, and the point $0$ is the smallest explicit witness.

Bounded primary-source searches through 17 August 2026 found the source but no
later arXiv paper recording this counterexample.  External expert review of
the ambient-boundary convention and Proposition 4.13(b) is recommended before
dissemination.

\begin{thebibliography}{9}
\raggedright
\bibitem{LN}
C.-W. Leung and C.-K. Ng,
\emph{Berkovich spectra of elements in Banach Rings},
arXiv:1410.5893 (2014; manuscript dated 2018).
\end{thebibliography}

\end{document}
